\section{Methodology}
\label{sec:method}

The central claim of this paper is that $\text{SE}_{N5}$ and min\_logprob are near-orthogonal uncertainty signals for open-weight LLMs on short-answer factual QA. To test this claim rigorously, we need: (1) precise signal definitions that make their algebraic distinction explicit, (2) a direct independence test protocol with pre-registered thresholds, and (3) a circularity-controlled evaluation design.

\subsection{Signal Definitions}

\textbf{Semantic Entropy ($\text{SE}_{N5}$).} Following \citet{kuhn2023semantic}, we define Semantic Entropy as:
\[
\text{SE} = H(p(C|x)) = -\sum_{C \in \mathcal{C}} p(C|x) \log p(C|x)
\]
where $p(C|x) = \sum_{s \in C} p(s|x)$ aggregates the token-probability mass of all $N=5$ stochastic samples within each NLI-defined semantic cluster $C$. Clusters are formed by bidirectional NLI entailment using \texttt{cross-encoder/nli-deberta-v3-small}. SE operates over $N=5$ stochastic samples (temperature$=0.7$, top-$p=0.95$), aggregating uncertainty at the semantic, sentence level. It does not depend on token-level probabilities from the greedy decode path.

\textbf{Minimum Log-Probability (min\_logprob).}
\[
\text{min\_logprob} = \min_{t \in [1,T]} \log p(w_t | w_{<t}, x)
\]
where the minimum is taken over all tokens in the greedy decode output of length $T$. min\_logprob captures the ``weakest link'' in the model's token-level confidence chain and is computed from a single greedy decoding pass \textbf{separate} from the $N=5$ stochastic sampling passes used for SE.

\textbf{Algebraic Distinctness.} The near-orthogonality claim follows from: (a) SE aggregates over the semantic clustering distribution of $N$ stochastic outputs; min\_logprob aggregates as the minimum over a greedy-path token sequence; (b) SE is computed on stochastic samples (temperature $> 0$); min\_logprob on the greedy decode (temperature $= 0$). The operations are not mathematically reducible to each other.

Figure~\ref{fig:heatmap} shows the full Pearson correlation matrix between $\text{SE}_{N5}$, min\_logprob, response length, and LM-judge correctness on $N=2500$ prompts --- confirming near-zero cross-correlation between SE and min\_logprob.

\subsection{Independence Test Protocol}

\textbf{Gate 1 --- Pearson Correlation Test.} Pre-registered threshold: PASS if $|r| < 0.70$; ABANDON if $|r| > 0.85$ (SE is a reparameterization of min\_logprob); EXPLORE if $0.70 \leq |r| \leq 0.85$.

\textbf{Gate 2 --- Conditional Logistic Regression (Partial $R^2$).} We fit:
\[
\text{logit}(h=1) \sim \beta_0 + \beta_1 \cdot \text{min\_logprob} + \beta_2 \cdot \text{SE} + \beta_3 \cdot L + \beta_4 \cdot (\text{SE} \times \text{min\_logprob})
\]
Partial $R^2(\text{SE})$ = McFadden $R^2$(full) $-$ McFadden $R^2$(reduced without SE), tested via LRT (2 df). PASS if partial $R^2 \geq 0.02$ AND LRT $p < 0.05$.

\subsection{Circularity Control Design}

A cross-model LM-judge (Qwen-2.5-7B-Instruct evaluating Llama-3.1-8B-Instruct outputs) prevents shared model-specific NLI biases from creating circular agreement between SE and judge labels. Circularity diagnostic: Spearman $\rho(\text{SE}, \text{judge\_correctness})$; threshold: $|\rho| < 0.40$.

\subsection{Implementation}

Six modules: \texttt{config.py} (hyperparameters), \texttt{generate.py} (checkpoint-aware stochastic + greedy generation), \texttt{compute\_signals.py} ($\text{SE}_{N5}$ + min\_logprob), \texttt{judge.py} (Qwen-2.5-7B LM-judge), \texttt{stats\_analysis.py} (Pearson $r$, conditional LR, partial $R^2$, gate evaluation), \texttt{visualize.py} (four figures). Checkpoint-aware pipeline saves signals to \texttt{results/signals.pkl}. GPU memory managed explicitly between generator and judge loads.
